\subsection{The case of any transcendental extension}

In this no. we shall use the following notation: K is a field, \(K'\) an extension of K, v a valuation on K, \(v'\) an extension of v to \(K'\) , \(\Gamma\) and k (resp. \(\Gamma'\) and \(k'\) ) the order group and residue field of \(v\) (resp. \(v'\) ). We shall write:

\[
d \left(\mathrm{K} ^ {\prime} / \mathrm{K}\right) = \dim . \mathrm{al} _ {\mathrm{K}} \mathrm{K} ^ {\prime} = \text { transcendence   degree   of } \mathrm{K} ^ {\prime} \text { over } \mathrm{K};
\]

\[
s (v ^ {\prime} / v) = \dim . \mathrm{al} _ {k} k ^ {\prime} = \text { transcendence   degree   of } k ^ {\prime} \text { over } k;
\]

\[
r (v ^ {\prime} / v) = r (\Gamma^ {\prime} / \Gamma) = \text { rational   rank   of } \Gamma^ {\prime} / \Gamma ,
\]

if the right-hand sides are finite; otherwise, we shall make the convention that \(d(\mathbf{K}' / \mathbf{K}) = +\infty\) (resp. \(s(v' / v) = +\infty\) , \(r(v' / v) = +\infty\) ).

THEOREM 1. Let \(x_1, \ldots, x_s\) be elements of the ring of \(v'\) whose canonical images \(\bar{x}_i\) in \(k'\) are algebraically independent over \(k\) and \(y_1, \ldots, y_r\) elements of \(K'\) such that the canonical images of the \(v'(y_j)\) in \(\Gamma'/\Gamma\) are linearly independent over \(Z\) . Then the \(r + s\) elements \(x_1, \ldots, x_s, y_1, \ldots, y_r\) of \(K'\) are algebraically independent over \(K\) ; the restriction of \(v'\) to \(K(x_1, \ldots, x_s, y_1, \ldots, y_r)\) admits \(k(\bar{x}_1, \ldots, \bar{x}_s)\) as residue field and

\[
\Gamma + \mathbf {Z} v ^ {\prime} (y _ {1}) + \dots + \mathbf {Z} v ^ {\prime} (y _ {r})
\]

as order group.

Our assertion is obvious if \(r + s = 0\) . We argue by induction on \(r + s\) . If \(r' \leqslant r\) , \(s' \leqslant s\) and \(r' + s' < r + s\) , the induction hypothesis shows that the hypotheses of Theorem 1 hold if \(K\) is replaced by \(K(x_1, \ldots, x_s, y_1, \ldots, y_{r'})\) and the families \((x_1, \ldots, x_s), (y_1, \ldots, y_r)\) by \((x_{s'+1}, \ldots, x_s), (y_{r'+1}, \ldots, y_r)\) . The problem is therefore reduced to one of the two following cases:

\begin{enumerate}
\def\labelenumi{(\alph{enumi})}
\item
  There is an element \(x\) in the ring of \(v'\) such that \(\bar{x}\) is transcendental over \(k\) ; then it is necessary to show that \(x\) is transcendental over \(\mathbf{K}\) and that the restriction of \(v'\) to \(\mathbf{K}(x)\) admits \(k(\bar{x})\) as residue field and \(\Gamma\) as order group.
\item
  There is an element \(y\) in \(\mathbf{K}'\) such that the relations \(nv'(y) \in \Gamma\) and \(n \in \mathbf{Z}\) imply \(n = 0\) ; it is necessary to show that \(y\) is transcendental over \(\mathbf{K}\) and that the restriction of \(v'\) to \(\mathbf{K}(y)\) admits \(k\) as residue field and \(\Gamma + \mathbf{Z}v'(y)\) as order group.
\end{enumerate}

Now Proposition 1 of §8, no. 1 shows that \(x\) (resp. \(y\) ) cannot be algebraic over \(K\) . The other assertions of (a) (resp. (b)) follow immediately by Proposition 2 (resp. Proposition 1) of no. 1.

COROLLARY 1. The inequality

\[
s (v ^ {\prime} / v) + r (v ^ {\prime} / v) \leqslant d (\mathbf {K} ^ {\prime} / \mathbf {K})\tag{9}
\]

holds.

Further, if \(K'\) is a finitely generated extension of K and equality holds in (9), then \(\Gamma'/\Gamma\) is a finitely generated Z-module and \(k'\) is a finitely generated extension of k.

Let \(r\) and \(s\) be natural numbers such that \(r \leqslant r(v'/v)\) and \(s \leqslant s(v'/v)\) ; we show that \(r + s \leqslant d(\mathbf{K}'/\mathbf{K})\) and this will prove (9). By hypothesis there exist elements \(x_1, \ldots, x_s, y_1, \ldots, y_r\) of \(\mathbf{K}'\) which satisfy the hypotheses of Theorem 1.

They are therefore algebraically independent over K, which shows the inequality \(r + s \leqslant d(\mathbf{K}'/\mathbf{K})\) .

If \(\mathbf{K}'\) is a finitely generated extension of \(\mathbf{K}\) , \(d(\mathbf{K}' / \mathbf{K})\) is finite, hence \(s(v' / v)\) and \(r(v' / v)\) are also finite; we denote them by \(s\) and \(r\) . There exist elements \(x_{1},\ldots ,x_{s},\) \(y_{1},\ldots ,y_{r}\) of \(\mathbf{K}'\) which satisfy the hypotheses of Theorem 1. If \(r + s = d(\mathbf{K}' / \mathbf{K})\) , these elements form a transcendence basis of \(\mathbf{K}'\) over \(\mathbf{K}\) and \(\mathbf{K}'\) is therefore a finite algebraic extension of \(\mathbf{K}'' = \mathbf{K}(x_1,\dots,y_r)\) . Let \(\Gamma''\) and \(k''\) be the order group and residue field of the restriction of \(v'\) to \(\mathbf{K}''\) . By Theorem 1, \(\Gamma'' / \Gamma\) is a finitely generated \(\mathbf{Z}\) -module and \(k''\) is a finitely generated pure extension of \(k\) . On the other hand, as \(\mathbf{K}'\) is a finite algebraic extension of \(\mathbf{K}''\) , \(\Gamma' / \Gamma\) is a finite group and \(k'\) is a finite algebraic extension of \(k''\) (§ 8, no. 1, Lemma 2). This proves the corollary.

COROLLARY 2. Let \(h\) and \(h'\) be the heights of \(v\) and \(v'\) . Then

\[
s (v ^ {\prime} / v) + h ^ {\prime} \leqslant d (\mathrm{K} ^ {\prime} / \mathrm{K}) + h.\tag{10}
\]

By Proposition 3, \(h' \leqslant r(v'/v) + h\) .

COROLLARY 3. Suppose that \(\mathbf{K}'\) is a finitely generated extension of \(\mathbf{K}\) , that \(\Gamma\) is isomorphic to \(\mathbf{Z}^h\) (ordered lexicographically) and there is equality in formula (10). Then \(\Gamma'\) is isomorphic to \(\mathbf{Z}^{h'}\) (ordered lexicographically) and \(k'\) is a finitely generated extension of \(k\) .

If there is equality in (10), there is equality in (9), whence the fact that \(k'\) is a finitely generated extension of k and that \(\Gamma'\) is a finitely generated Z-module. Further, comparing (9) and (10), it is seen that \(h' - h = r(\Gamma'/\Gamma)\) , whence \(h' = r(\Gamma')\) and Proposition 4 (no. 2) then shows that \(\Gamma'\) is isomorphic to \(Z^{h'}\) ordered lexicographically.

COROLLARY 4. Suppose that \(v\) is improper (in which case \(k = K\) ). Then

\[
h (\Gamma^ {\prime}) + d (k ^ {\prime} / \mathbf {K}) \leqslant r (\Gamma^ {\prime}) + \mathbf {D} (k ^ {\prime} / \mathbf {K}) \leqslant d (\mathbf {K} ^ {\prime} / \mathbf {K}).\tag{11}
\]

If, in particular, \(v'\) is of height 1, then

\[
d (k ^ {\prime} / \mathbf {K}) \leqslant d (\mathbf {K} ^ {\prime} / \mathbf {K}) - 1;\tag{12}
\]

further, if \(\mathbf{K}'\) is a finitely generated extension of \(\mathbf{K}\) and there is equality in (12), then \(v'\) is a discrete valuation and \(k'\) is a finitely generated extension of \(\mathbf{K}\) .

It is a series of special cases of Corollaries 1, 2, 3.


